% Exemple complet utilisant toutes les fonctionnalités du package

% === Barème avec badges ===
\bareme{
	\exbadge{1}{4} \quad
	\exbadge{2}{6} \quad
	\exbadge{3}{5} \quad
	\exbadge{4}{7} \quad
	\bonusbadge{2}
}

% === Consignes ===
\begin{consignebox}
	\faIcon{info-circle} \textbf{CONSIGNES IMPORTANTES}
	\begin{itemize}[leftmargin=*,noitemsep,topsep=5pt]
		\item[\faIcon{check}] Répondez directement sur le sujet
		\item[\faIcon{check}] Justifiez toutes vos réponses
		\item[\faIcon{pencil-alt}] Soignez la présentation
		\item[\faIcon{times}] Calculatrice non autorisée
	\end{itemize}
\end{consignebox}

\exspace

% === Grille de compétences (optionnel) ===
\CompetencesGrid{
	\item[\faIcon{check}] Chercher : Extraire les informations utiles d'un énoncé
	\item[\faIcon{check}] Modéliser : Traduire une situation en langage mathématique
	\item[\faIcon{check}] Calculer : Effectuer des calculs et vérifier les résultats
	\item[\faIcon{check}] Communiquer : Rédiger une solution claire et organisée
}

\exspace

% ============================================
% EXEMPLE 1 : QCM avec la nouvelle boîte
% ============================================
\begin{qcmbox}{\faIcon{list-ul} QCM : Nombres entiers \points{4}}
	Pour chaque question, écris la lettre de la seule réponse correcte dans la colonne de droite :
	
	\vspace{0.3cm}
	
	\begin{qcmtable}
		\textbf{1.} Combien vaut $5 + 7$ ? & 11 & 12 & 13 & \ans{B} \\
		\hline
		\textbf{2.} Quel est le double de 15 ? & 25 & 30 & 35 & \ans{B} \\
		\hline
		\textbf{3.} $100 - 47 = $ ? & 43 & 53 & 63 & \ans{B} \\
		\hline
		\textbf{4.} Combien de dizaines dans 235 ? & 2 & 3 & 23 & \ans{C} \\
		\hline
	\end{qcmtable}
\end{qcmbox}

\exspace

% ============================================
% EXEMPLE 2 : Exercice avec template rapide
% ============================================
\ExTemplate{2}{Opérations posées}{6}{calculator}{

\textbf{Effectue les opérations suivantes en les posant :}

\smallspace

\textbf{1.} $345 + 278$ \points{1.5}

\begin{center}
	\AddPose{345}{278}
\end{center}

\exspace

\textbf{2.} $534 - 187$ \points{1.5}

\begin{center}
	\SubPose{534}{187}
\end{center}

\exspace

\textbf{3.} $48 \times 23$ \points{1.5}

\begin{center}
	\MulPose{48}{23}
\end{center}

\exspace

\textbf{4.} $156 \div 12$ (division euclidienne) \points{1.5}

\begin{center}
	\DivEuc{156}{12}
\end{center}

Donc : $156 = 12 \times$ \ans{13} $+$ \ans{0}

}

% ============================================
% EXEMPLE 3 : Construction géométrique
% ============================================
\begin{constructionbox}{\faIcon{drafting-compass} Exercice 3 : Construction \points{5}}

\textbf{Construis la figure suivante avec tes instruments de géométrie :}

\begin{itemize}[leftmargin=*]
	\item Un segment $[AB]$ de longueur 6 cm
	\item Un point $C$ tel que $AC = 4$ cm et $BC = 5$ cm
	\item La perpendiculaire à $(AB)$ passant par $C$
	\item Le cercle de centre $A$ et de rayon 3 cm
\end{itemize}

\figureespace{8cm}

\onlycorrige{
	\textcolor{warning}{\textbf{Points de notation :}}
	\begin{itemize}[leftmargin=*]
		\item Segment [AB] correctement tracé : 1 pt
		\item Point C bien placé : 1.5 pts
		\item Perpendiculaire tracée et codée : 1.5 pts
		\item Cercle tracé : 1 pt
	\end{itemize}
}

\end{constructionbox}

\exspace

% ============================================
% EXEMPLE 4 : Problème avec zones de réponse
% ============================================
\begin{problembox}{\faIcon{search} Exercice 4 : Problème de vitesse \points{7}}

\begin{tcolorbox}[colback=blue!5,colframe=primary!30,arc=2mm,boxrule=0.5pt]
	\textbf{Situation :} Un train parcourt 180 km en 2 heures.
\end{tcolorbox}

\exspace

\textbf{1.} Quelle est la vitesse moyenne du train en km/h ? \points{2}

\anslines[0.5cm]{3}{
Vitesse = Distance $\div$ Temps

Vitesse = $180 \div 2 = 90$ km/h

La vitesse moyenne du train est de 90 km/h.
}

\exspace

\textbf{2.} Quelle distance parcourt-il en 3 heures à cette vitesse ? \points{2}

\anslines[0.5cm]{3}{
Distance = Vitesse $\times$ Temps

Distance = $90 \times 3 = 270$ km

Le train parcourt 270 km en 3 heures.
}

\exspace

\textbf{3.} Combien de temps met-il pour parcourir 450 km ? \points{3}

\anslines[0.5cm]{4}{
Temps = Distance $\div$ Vitesse

Temps = $450 \div 90 = 5$ heures

Le train met 5 heures pour parcourir 450 km.
}

\end{problembox}

\exspace

% ============================================
% BONUS : Avec opération posée et raisonnement
% ============================================
\begin{tcolorbox}[
	colback=accent!10,
	colframe=accent,
	fonttitle=\bfseries,
	title={\faIcon{star} BONUS : Défi mathématique \bonusbadge{2}},
	arc=3mm,
	boxrule=1.5pt,
	enhanced,
	width=\linewidth,
	grow to left by=5mm,
	grow to right by=5mm
]

Un nombre mystérieux vérifie les propriétés suivantes :
\begin{itemize}[leftmargin=*]
	\item Il est compris entre 100 et 200
	\item La somme de ses chiffres vaut 15
	\item Il est divisible par 3 et par 5
\end{itemize}

\textbf{Quel est ce nombre ? Justifie ta réponse.}

\anslines[0.5cm]{8}{
Le nombre doit être divisible par 3 ET par 5, donc divisible par 15.

Multiples de 15 entre 100 et 200 : 105, 120, 135, 150, 165, 180, 195

Vérification de la somme des chiffres = 15 :

• 105 : $1+0+5 = 6$ ✗ \quad • 120 : $1+2+0 = 3$ ✗ \quad • 135 : $1+3+5 = 9$ ✗

• 150 : $1+5+0 = 6$ ✗ \quad • 165 : $1+6+5 = 12$ ✗ \quad • 180 : $1+8+0 = 9$ ✗

• 195 : $1+9+5 = 15$ ✓

Le nombre mystérieux est \textbf{195}.
}

\end{tcolorbox}

\exspace

% ============================================
% DÉMONSTRATION : Utilisation des helpers
% ============================================
\onlycorrige{
\begin{tcolorbox}[colback=warning!10,colframe=warning,
	title={\faIcon{lightbulb} Notes pour le professeur},
	arc=2mm]
	
	\textbf{Barème total : 24 points}
	
	\begin{itemize}[leftmargin=*]
		\item QCM : 4 pts (1 pt par question)
		\item Opérations : 6 pts (1.5 pt par opération)
		\item Construction : 5 pts (voir détail dans l'exercice)
		\item Problème : 7 pts (2+2+3)
		\item Bonus : 2 pts
	\end{itemize}
	
	\textbf{Note finale sur 20 :} $(points \times 20) \div 24$
	
	\textbf{Temps indicatif :}
	\begin{itemize}
		\item QCM : 8 min
		\item Opérations : 15 min
		\item Construction : 12 min
		\item Problème : 15 min
		\item Bonus : 5 min
	\end{itemize}
\end{tcolorbox}
}

% Auto-évaluation
\AutoEval
